\section{MonPD Reborn}
\label{sec:monpd-reborn}

MonPD Reborn adalah inisiatif pengembangan ulang (redevelopment) dari sistem MonPD terdahulu, yang bertujuan untuk memodernisasi arsitektur sistem dan mengatasi seluruh kendala performa yang ada. Proyek ini bukan sekadar pembaruan tampilan, melainkan pembangunan ulang inti aplikasi menggunakan teknologi terkini berbasis .NET dengan framework ASP.NET Core MVC.

MonPD Reborn dirancang dengan fokus utama pada efisiensi performa dan pengalaman pengguna (User Experience). Perbedaan mendasar dan pembaharuan utama yang ditawarkan oleh MonPD Reborn dibandingkan pendahulunya meliputi:

\begin{enumerate}
    \item Penerapan Teknologi Asinkron (AJAX): \\
    Sistem baru ini mengubah paradigma pemuatan data dari sinkron menjadi asinkron. Aplikasi hanya memuat kerangka halaman terlebih dahulu, kemudian data diambil secara terpisah di latar belakang sesuai permintaan pengguna. Hal ini membuat aplikasi terasa jauh lebih responsif dan ringan.
    \item Antarmuka Interaktif (DevExtreme): \\
    MonPD Reborn mengimplementasikan komponen antarmuka modern yang memungkinkan pengguna melakukan interaksi data tingkat lanjut secara langsung, seperti pemfilteran per kolom, pengurutan data (sorting), dan navigasi halaman (paging) di sisi server, tanpa perlu memuat ulang halaman.
    \item Fitur Visualisasi Data: \\ 
    Sistem ini dilengkapi dengan dashboard analitik yang menyajikan data dalam bentuk grafik, card KPI (Key Performance Indicator), dan pemetaan geospasial (peta sebaran objek pajak) untuk memudahkan pegawai dalam membaca tren pendapatan daerah.
    \item Modul Komprehensif: \\
    MonPD Reborn mengintegrasikan berbagai modul pengawasan yang mencakup Monitoring Objek Pajak, Pengawasan Reklame, Aktivitas Penagihan Wajib Pajak, hingga Laporan Strategis dalam satu platform terpadu.
\end{enumerate}

Melalui pengembangan MonPD Reborn, Bapenda Kota Surabaya menargetkan terciptanya ekosistem kerja digital yang lebih cepat, akurat, dan efisien dalam rangka optimalisasi Pendapatan Asli Daerah (PAD). 